%% ma: L12-B16-0016
%% lop: 12
%% chuong: 5
%% bai: 16
%% dang: 12.16.9
%% loai: TN4
%% muc: TH
%% dap_an: "B"
%% nguon: "Ảnh giáo viên gửi 04/11/2026 (Huỳnh Văn Ánh, Chuyên đề V - Phương pháp tọa độ trong không gian), Đề test theo bài, Phần I Câu 7"
%% kien_thuc: "Bài 16 (góc giữa hai mặt phẳng), Bài 14 (vectơ pháp tuyến, phương trình mặt phẳng tọa độ)"
%% trang_thai: chinh-thuc
%% ghi_chu: "Đáp án gốc B khớp (sympy). Đã duyệt 04/11/2026."
\begin{ex}
Trong không gian với hệ tọa độ $Oxyz$, tính góc giữa mặt phẳng $(P)\colon 2x-3y-4=0$ và mặt phẳng $(Oxy)$.
\choice{$30^\circ$}{\True $90^\circ$}{$60^\circ$}{$0^\circ$}
\loigiai{$(P)$ có vectơ pháp tuyến $\vec n_1=(2;-3;0)$. Mặt phẳng $(Oxy)$ có phương trình $z=0$ nên có vectơ pháp tuyến $\vec n_2=(0;0;1)$.

Do đó $\cos((P),(Oxy))=\dfrac{|\vec n_1\cdot\vec n_2|}{|\vec n_1|\cdot|\vec n_2|}=0$, suy ra $((P),(Oxy))=90^\circ$. Chọn B.}
\end{ex}
